\begin{acknowledgments}
The author used AI assistants in this work. Anthropic Claude was used in two separate sessions. The first session (claude-opus-5-5) wrote the instructions for each stage, re-ran every verification suite in a fresh folder, recomputed the key constants by hand, and reviewed the manuscript. The second session derived the lemmas and the proofs of Theorems 2 to 5, wrote the interval arithmetic and the verification code, and drafted the text, using Anthropic Claude Fable (claude-fable-5-2) for the new proof ideas, among them the narrow-focus lemmas and the positivity-and-periodicity lemma (Lemma 7), and claude-opus-5-5 otherwise; it also found constants that had been left as numerical maxima and closed them by proofs. Google Gemini (Gemini 4 Argon) and OpenAI ChatGPT (GPT-6 Sol) independently reviewed the manuscript; their comments led to the lower end of the borrowed Lemma BE and several corrections of notation.
The author directed the research, made all decisions on its content, and takes full responsibility for this article.
\end{acknowledgments}

\emph{Data and code availability.}
All code, stored results, the interval-arithmetic computations and the verification suite are available at \url{https://github.com/neginyan/resolution-gap-every-focus} and archived at Zenodo (DOI added on publication).
A single entry point, \code{verification/run\_all.py}, recomputes every computer-assisted part of this paper and re-derives its numbers from the stored results; \code{verification/quick.py} is a short check that runs in a few minutes.
The certificates of Theorem~H1 are those of~\cite{H} and are re-run with the repository of~\cite{H}.
Paper~F and its repository~\cite{F} are not changed; this paper uses version~1.3 of paper~H~\cite{H}, prepared together with it.
